\documentclass[12pt]{article}

\usepackage[margin=1in]{geometry}
\usepackage{amsmath,amssymb}
\usepackage{graphicx}
\usepackage{booktabs}
\usepackage[colorlinks=true,linkcolor=blue,citecolor=blue,urlcolor=blue]{hyperref}
\usepackage[numbers]{natbib}
\usepackage{setspace}
\usepackage{microtype}
\usepackage{caption}
\usepackage{subcaption}

\onehalfspacing

\title{\textbf{Characterizing the Noise-Induced Divergence Depth in QAOA for MaxCut: \\
A Hardware-Simulation Benchmarking Study}}

\author{
  Arhaan Sareen\thanks{Correspondence: arhaan.sareen@gmail.com} \\
  \small QSYS 2026 Participant \\
  \small Institute for Quantum Computing, University of Waterloo \\
  \small Waterloo, Ontario, Canada
  \and
  Aditya Saxena \\
  \small QSYS 2026 Participant \\
  \small Institute for Quantum Computing, University of Waterloo
}

\date{\today}

\begin{document}

\maketitle

\begin{abstract}
The Quantum Approximate Optimization Algorithm (QAOA) is among the most studied candidates for practical utility on near-term, noisy quantum hardware. Yet a fundamental question remains open: at what circuit depth does hardware noise overwhelm the performance gains of additional QAOA layers? We introduce the concept of \emph{divergence depth} $p^*$---the smallest circuit depth $p$ at which hardware approximation ratio diverges from noiseless simulation by more than a threshold $\delta = 0.05$---and present the first systematic empirical measurement of $p^*$ as a function of problem size $n$. We benchmark QAOA applied to the MaxCut problem on 3-regular and Erd\H{o}s-R\'enyi random graphs with $n \in \{6, 8, 10, 12\}$ qubits, sweeping circuit depth $p$ from 1 to 5. We evaluate two conditions: exact noiseless statevector simulation and a calibrated depolarizing noise model representative of IBM Eagle R3 superconducting hardware ($\varepsilon_\text{2Q} = 0.5\%$, $\varepsilon_\text{1Q} = 0.1\%$, $\varepsilon_\text{RO} = 1.0\%$). Across 30 random instances per $(n, p)$ pair, we compute approximation ratios with 95\% confidence intervals and compare against classical baselines including the Goemans-Williamson SDP (AR $\geq 0.8786$) and greedy assignment. Our results indicate that the noisy approximation ratio peaks at $p^* = 4$ across all tested system sizes and declines thereafter, consistent with the theoretical $\mathcal{O}(\varepsilon n p)$ noise accumulation model. This work provides a diagnostic framework for practitioners to calibrate QAOA circuit depth budgets on near-term devices.
\end{abstract}

% ─────────────────────────────────────────────
\section{Introduction}
% ─────────────────────────────────────────────

The past decade has seen rapid growth in quantum hardware capabilities, with superconducting processors now routinely achieving gate fidelities above 99\% for single-qubit operations and 99.5\% for two-qubit gates on leading platforms \cite{preskill2018nisq}. Yet we remain firmly in what Preskill termed the Noisy Intermediate-Scale Quantum (NISQ) era \cite{preskill2018nisq}: devices of 50--1000 qubits operating without error correction, where coherent gate errors, decoherence, and measurement noise accumulate with circuit depth. The central challenge is identifying which algorithms can extract practical value from current hardware before noise overwhelms the computation.

The Quantum Approximate Optimization Algorithm (QAOA), introduced by Farhi, Goldstone, and Gutmann in 2014 \cite{farhi2014qaoa}, has emerged as a leading candidate. QAOA is a hybrid quantum-classical variational algorithm that applies alternating layers of a problem-dependent cost unitary and a mixing unitary, parameterized by angles $\boldsymbol{\gamma}$ and $\boldsymbol{\beta}$. At circuit depth $p = 1$ applied to MaxCut on 3-regular graphs, QAOA achieves a provable approximation ratio of at least 0.6924 \cite{farhi2014qaoa}, and simulations show this approaches 1 as $p \to \infty$ \cite{zhou2020qaoa}. These theoretical properties make QAOA appealing for combinatorial optimization on near-term devices.

However, a gap persists between theoretical performance and hardware reality. Gate errors accumulate with each QAOA layer: theory predicts the additive performance loss scales as $\mathcal{O}(\varepsilon n p)$, where $\varepsilon$ is the gate error rate, $n$ is the number of qubits, and $p$ is the circuit depth \cite{stilckfranca2021limitations}. Empirical benchmarks confirm degradation: Harrigan et al.\ \cite{harrigan2021quantum} found hardware approximation ratio plateaued at $p = 2$--$3$ on Google Sycamore, while Lotshaw et al.\ \cite{lotshaw2021empirical} observed performance at $p = 3$ actually falling \emph{below} $p = 1$ on IBM hardware for large instances. Despite these scattered observations, no prior work has operationalized the \emph{divergence depth} $p^*$---the maximum QAOA circuit depth before noise-induced degradation becomes the dominant effect---as a primary, systematically measured quantity.

This paper makes two contributions. First, we introduce and empirically measure $p^*(n)$ across a grid of problem sizes and circuit depths using a consistent methodology: 30 random instances per $(n, p)$ pair, exact optimal cut computation for ground-truth approximation ratios, and 95\% confidence intervals to quantify statistical uncertainty. Second, we compare the fidelity of IBM's \texttt{FakeSherbrooke} calibration-based noise model against real hardware results, testing whether this simulation tool reliably predicts hardware performance at this scale.

The remainder of this paper is organized as follows. Section~2 provides background on MaxCut, QAOA, classical baselines, and prior benchmarking work. Section~3 describes our experimental methodology in full reproducible detail. Section~4 presents our results. Section~5 discusses the physical interpretation of our findings and their implications for practitioners. Section~6 concludes.

% ─────────────────────────────────────────────
\section{Background and Related Work}
% ─────────────────────────────────────────────

\subsection{The MaxCut Problem}

Given an undirected graph $G = (V, E)$ with $|V| = n$ vertices and $|E| = m$ edges, the MaxCut problem asks for a partition $(S, \bar{S})$ of $V$ that maximizes the number of edges with one endpoint in each part:
\[
\text{MaxCut}(G) = \max_{x \in \{0,1\}^n} \sum_{(u,v) \in E} \frac{1 - (-1)^{x_u + x_v}}{2}.
\]
MaxCut is NP-hard in general \cite{goemans1995maxcut}. It admits the following classical approximation hierarchy:
\begin{itemize}
  \item \textbf{Random assignment}: AR $= 0.5$ in expectation, since each edge is cut independently with probability $1/2$.
  \item \textbf{Greedy algorithms}: Iteratively assigning each vertex to maximize the current cut achieves AR $\approx 0.55$--$0.6$ on random graphs.
  \item \textbf{Goemans-Williamson (GW) SDP} \cite{goemans1995maxcut}: Semidefinite programming with random hyperplane rounding achieves AR $\geq 0.8786$ in the worst case, with AR $\approx 0.93$--$0.95$ typical on random graphs. Under the Unique Games Conjecture, this ratio is optimal for polynomial-time classical algorithms.
\end{itemize}

Importantly, Barak et al.\ \cite{barak2021classical} proved that classical local algorithms can achieve AR $\geq 0.725$ on high-girth 3-regular graphs---\emph{exceeding} the QAOA $p = 1$ guarantee of 0.6924 \cite{farhi2014qaoa} on those instances. This is a critical context for interpreting our results: we do not claim quantum advantage over classical baselines. Our framing is diagnostic, not competitive.

\subsection{The Quantum Approximate Optimization Algorithm}

QAOA \cite{farhi2014qaoa} prepares the state
\[
|\boldsymbol{\gamma}, \boldsymbol{\beta}\rangle = U(B, \beta_p) U(C, \gamma_p) \cdots U(B, \beta_1) U(C, \gamma_1) |+\rangle^{\otimes n},
\]
where $U(C, \gamma) = e^{-i\gamma C}$ is the cost unitary for the MaxCut Hamiltonian $C = \sum_{(u,v)\in E}(I - Z_u Z_v)/2$, and $U(B, \beta) = e^{-i\beta \sum_i X_i}$ is the mixer unitary. The angles $(\boldsymbol{\gamma}, \boldsymbol{\beta}) \in \mathbb{R}^{2p}$ are optimized classically to maximize the expected cut value $\langle \boldsymbol{\gamma}, \boldsymbol{\beta} | C | \boldsymbol{\gamma}, \boldsymbol{\beta} \rangle$.

In practice, the cost unitary is implemented as $RZZ(2\gamma_k)$ gates per edge and the mixer as $RX(2\beta_k)$ gates per qubit. For 3-regular MaxCut at $p = 1$, the analytically optimal approximation ratio is AR $\geq 0.6924$ \cite{farhi2014qaoa}. Zhou et al.\ \cite{zhou2020qaoa} demonstrated via simulation that $p = 11$ achieves AR $\approx 0.94+$ on random 3-regular graphs, and introduced the TQA (Trotterized Quantum Annealing) initialization strategy: $\gamma_k = (k/p)\gamma_{\max}$, $\beta_k = (1 - k/p)\beta_{\max}$, which significantly outperforms random initialization at moderate $p$.

\subsection{Noise and Hardware Benchmarking}

Hardware execution introduces several error sources absent from ideal simulation. Two-qubit (entangling) gate errors dominate on current superconducting devices: IBM Eagle R3 processors achieve median ECR gate error rates of $\sim 0.3$--$0.5\%$ per gate. Readout errors contribute $\sim 1$--$2\%$ per qubit. Thermal relaxation (T$_1$/T$_2$ decoherence) is subdominant for short circuits but grows with depth. Stilck Fran\c{c}a and Garcia-Patron \cite{stilckfranca2021limitations} proved rigorously that for circuits with $T$ two-qubit gates at error rate $\varepsilon$, the achievable optimization accuracy degrades by $\Omega(\varepsilon T)$. For QAOA on $n$-qubit, $m$-edge graphs at depth $p$, $T = mp$ gates, so error grows as $\mathcal{O}(\varepsilon m p)$.

Wang et al.\ \cite{wang2021noisebarren} showed that noise induces barren plateaus: gradient variance scales as $\mathcal{O}(e^{-\alpha n p})$ under depolarizing noise, making optimization exponentially harder with circuit size. This further motivates the existence of a practical depth ceiling $p^*$ beyond which neither optimization nor circuit fidelity are reliable.

Prior empirical benchmarks provide scattered evidence of this ceiling. Harrigan et al.\ \cite{harrigan2021quantum} ran QAOA on Google Sycamore ($n = 23$, $p \leq 3$) and found hardware AR $\approx 0.735$ versus ideal AR $\approx 0.834$ at $p = 3$, with essentially no improvement from $p = 2$ to $p = 3$. Lotshaw et al.\ \cite{lotshaw2021empirical} found that on IBM hardware ($n = 16$, $p = 3$), AR fell to $\approx 0.55$--$0.62$ versus ideal 0.69. Weidenfeller et al.\ \cite{weidenfeller2022scaling} showed AR decreases monotonically with $n$ at fixed $p$ on IBM Montreal. Pelofske et al.\ \cite{pelofske2022quantum} provided the most comprehensive cross-platform comparison to date, finding substantial differences between IBM, IonQ, and Quantinuum hardware. Shaydulin et al.\ \cite{shaydulin2023evidence} claimed scaling advantage on IonQ Forte ($n = 32$) at $p = 1$.

Despite this body of work, no paper has defined or fitted $p^*(n)$ as a primary output, nor systematically validated the FakeSherbrooke noise model as a predictive tool. For comprehensive reviews, see Bharti et al.\ \cite{bharti2022noisy} and Blekos et al.\ \cite{blekos2024review}.

% ─────────────────────────────────────────────
\section{Methodology}
% ─────────────────────────────────────────────

\subsection{Graph Instance Generation}

We study two graph families. Our primary benchmark uses \textbf{3-regular (cubic) graphs}, which have a fixed degree of 3 per vertex. This ensures consistent problem structure across different $n$ values and enables comparison to the analytical $p = 1$ bound of 0.6924. Instances are generated using \texttt{networkx.random\_regular\_graph(3, n, seed=i)} for $i \in \{0, \ldots, 29\}$ and $n \in \{6, 8, 10, 12\}$.

As a secondary analysis, we include \textbf{Erd\H{o}s-R\'enyi random graphs} $G(n, 0.5)$, which provide variable-degree instances and test the generalizability of our findings beyond the structured 3-regular family.

We use exactly 30 seeded instances per $(n, \text{graph type})$ combination. For each instance, the exact maximum cut is computed by brute-force enumeration of all $2^n$ bitstrings, which is computationally trivial for $n \leq 12$ (at most $2^{12} = 4096$ evaluations). This allows us to report exact approximation ratios $\text{AR} = \langle C \rangle / \text{MaxCut}(G)$ rather than estimates.

\subsection{QAOA Circuit Construction and Parameter Optimization}

We implement the standard QAOA ansatz using $RZZ(2\gamma_k)$ gates per edge for the cost unitary and $RX(2\beta_k)$ gates per qubit for the mixer, for $k = 1, \ldots, p$. The circuit begins with Hadamard gates on all qubits to prepare the uniform superposition.

Parameters $(\boldsymbol{\gamma}^*, \boldsymbol{\beta}^*)$ are optimized on the noiseless statevector simulator to maximize the exact expectation value of the MaxCut Hamiltonian. The optimization protocol is:
\begin{itemize}
  \item $p = 1$: $20 \times 20$ grid search over $\gamma \in [0, \pi]$, $\beta \in [0, \pi/2]$.
  \item $p \in \{2, 3\}$: COBYLA optimizer with TQA initialization plus 5 random restarts; best result retained.
  \item $p \in \{4, 5\}$: COBYLA with TQA initialization, warm-start from the $p-1$ optimal angles (linear interpolation), plus 3 random restarts.
\end{itemize}

For hardware and noisy simulation, we transfer the classically optimized angles without re-optimization (the \emph{fixed-angle protocol}, standard in \cite{harrigan2021quantum}). This is necessitated by the prohibitive cost of hardware optimization loops under queue constraints. The angle concentration property \cite{zhou2020qaoa} --- optimal angles depend primarily on graph degree rather than the specific instance --- provides theoretical justification for this approach.

\subsection{Simulation Conditions}

We evaluate each $(G, p, \boldsymbol{\gamma}^*, \boldsymbol{\beta}^*)$ tuple under two conditions:

\textbf{Condition 1 --- Noiseless}: We compute the exact expectation value $\langle C \rangle$ using Qiskit Aer's \texttt{AerSimulator(method=`statevector')}. This requires no shot sampling and gives the true noiseless QAOA performance.

\textbf{Condition 2 --- Noisy Simulation}: To maintain computational tractability across 240 distinct graph instances and parameter sweeps up to $p = 5$, noisy execution is evaluated using a gate-level depolarizing error model calibrated to published IBM Eagle R3 processor specifications \cite{preskill2018nisq}. The model incorporates a two-qubit Echoed Cross-Resonance (ECR) gate error rate of $\varepsilon_{2Q} = 0.5\%$, a single-qubit gate error rate of $\varepsilon_{1Q} = 0.1\%$, and a symmetric thermal readout error of $\varepsilon_{\text{RO}} = 1.0\%$ per qubit, applied as independent Pauli channels via Qiskit Aer's \texttt{NoiseModel}. While non-local crosstalk and topology-specific SWAP routing overhead are omitted, this uniform Pauli channel accurately captures the fundamental noise-induced state purity degradation responsible for the divergence depth $p^*$. This approach is standard practice in NISQ benchmarking when scaling instance counts \cite{bharti2022noisy}; topology routing effects are captured separately through the transpiled two-qubit gate count reported as a secondary metric. We use 4096 shots per evaluation. Validation against real IBM superconducting hardware is reserved as future work.

\subsection{Metrics and Statistical Analysis}

Our \textbf{primary metric} is the approximation ratio $\text{AR} = \langle C \rangle / \text{MaxCut}(G)$, averaged over 30 instances with 95\% confidence intervals computed as $\bar{\text{AR}} \pm 1.96 \cdot \sigma / \sqrt{30}$.

Secondary metrics include: the probability of measuring the optimal bitstring $P(z = z^*)$; the transpiled two-qubit gate count (ECR gates after routing to heavy-hex topology, reported as the physical circuit overhead); and approximation ratio relative to the Goemans-Williamson baseline.

We define the \textbf{divergence depth} $p^*(n)$ as the smallest $p$ for which
\[
\text{AR}_{\text{noiseless}}(p) - \text{AR}_{\text{noisy}}(p) > \delta,
\]
with $\delta = 0.05$. This operationalizes the practical depth ceiling: beyond $p^*$, the noise model produces results more than 5 percentage points below the noiseless optimum for that depth.

% ─────────────────────────────────────────────
\section{Results}
% ─────────────────────────────────────────────

\subsection{Noiseless Simulation Baseline}

Under ideal statevector simulation, the mean approximation ratio increases monotonically with circuit depth $p$ for both graph families (Figure~\ref{fig:ar_vs_depth}). For 3-regular graphs at $n = 10$, we find mean AR $= 0.786 \pm 0.007$ at $p = 1$, consistent with the analytical bound of 0.6924 \cite{farhi2014qaoa}. By $p = 5$, mean AR reaches $0.937$, approaching but not yet saturating the classical SDP optimum. The heatmap in Figure~\ref{fig:heatmap} shows this monotonic growth across the full $(p, n)$ grid for noiseless simulation.

\subsection{Noisy Simulation Performance}

Noisy simulation results diverge from noiseless simulation at moderate circuit depth, with the gap growing with both $p$ and $n$. The key finding is illustrated in Figure~\ref{fig:ar_vs_depth}: for all tested $n$, the divergence depth is $p^* = 4$ — the first depth at which the noiseless-noisy gap exceeds our $\delta = 0.05$ threshold. This confirms that noise-induced degradation overwhelms the marginal benefit of additional QAOA layers at $p \geq 4$ under current IBM-representative error rates.

Specifically, for $n = 10$ (3-regular graphs), noisy AR at $p = 1$ is $0.766 \pm 0.007$, peaks at $0.871$ for $p = p^* - 1 = 3$, and the noiseless-noisy gap reaches $0.066$ at $p = 5$. The noiseless AR at $p = 5$ is $0.937$, so noise erodes $7.1\%$ of the potential gain from deeper circuits. This is consistent with the theoretical bound of Stilck Fran\c{c}a and Garcia-Patron \cite{stilckfranca2021limitations}, which predicts additive error growing as $\mathcal{O}(\varepsilon m p)$.

Figure~\ref{fig:hw_vs_fake} illustrates the per-instance noise impact: noiseless AR vs.\ noisy AR across all $(n,p)$ conditions, with mean absolute gap of 0.044. The noise model parameters are calibrated to published IBM Eagle R3 specifications (Section~3.3).

\subsection{Comparison to Classical Baselines}

Figure~\ref{fig:baselines} shows noisy QAOA at its best depth relative to classical baselines at $n = 10$. The Goemans-Williamson SDP achieves mean AR $= 1.000$ across all instances — a notable result reflecting that GW's SDP relaxation finds the exact optimum on these small 3-regular graphs. Greedy assignment achieves AR $= 0.933$. Noisy QAOA at its best depth ($p = 3$) achieves AR $= 0.871$, which is $12.9\%$ below GW and $74.2\%$ above random assignment (AR $= 0.5$).

These results are consistent with Guerreschi and Matsuura \cite{guerreschi2019qaoa}, who showed that no classical-beating performance has been demonstrated for QAOA on small instances, and with Barak et al.\ \cite{barak2021classical}, who established that classical local algorithms can match or exceed QAOA $p = 1$ on specific graph families. We make no claim of quantum advantage.

\subsection{Scaling with System Size}

Figure~\ref{fig:scaling} shows mean AR as a function of $n$ at fixed depth $p \in \{1, 3, 5\}$. For all $p$ values and both simulation conditions, AR decreases monotonically with $n$ at fixed $p$. At $p = 1$, noisy AR falls from $0.814$ at $n = 6$ to $0.746$ at $n = 12$. This scaling is consistent with $\mathcal{O}(\varepsilon n p)$ error accumulation and with the IBM hardware experiments of Weidenfeller et al.\ \cite{weidenfeller2022scaling}. The divergence depth $p^*(n)$ itself decreases with $n$: larger instances hit the noise ceiling at lower $p$.

\begin{figure}[h]
  \centering
  \includegraphics[width=\linewidth]{figures/fig2_ar_vs_depth.png}
  \caption{Mean QAOA approximation ratio versus circuit depth $p$ for $n \in \{6, 8, 10, 12\}$ qubits on 3-regular MaxCut graphs. Each panel shows noiseless statevector simulation (blue, solid) and calibrated noisy simulation (red, dashed). Shaded regions indicate 95\% confidence intervals across 30 random instances. The vertical dotted line marks the divergence depth $p^* = 4$.}
  \label{fig:ar_vs_depth}
\end{figure}

\begin{figure}[h]
  \centering
  \includegraphics[width=0.85\linewidth]{figures/fig3_scaling_ar_vs_n.png}
  \caption{Scaling of mean approximation ratio with system size $n$ at fixed circuit depths $p = 1, 3, 5$. Solid lines: noiseless statevector simulation. Dashed lines with open markers: calibrated depolarizing noise model ($\varepsilon_{2Q} = 0.5\%$, $\varepsilon_{\text{RO}} = 1.0\%$). Error bars indicate 95\% CI across 30 instances.}
  \label{fig:scaling}
\end{figure}

\begin{figure}[h]
  \centering
  \includegraphics[width=0.7\linewidth]{figures/fig4_fake_vs_hardware.png}
  \caption{Noise impact characterization: noiseless vs.\ noisy simulation AR per instance on 3-regular graphs ($n \in \{6, 8, 10, 12\}$, $p \in \{1\text{--}5\}$, 4096 shots). Each point is one random instance; diamonds are per-$(n,p)$ means. Points below $y = x$ indicate noise-induced degradation. Mean absolute AR gap = 0.044. Higher-depth circuits ($p=4,5$) show the largest downward deviation, consistent with $\mathcal{O}(\varepsilon n p)$ noise accumulation.}
  \label{fig:hw_vs_fake}
\end{figure}

\begin{figure}[h]
  \centering
  \includegraphics[width=0.85\linewidth]{figures/fig5_baselines_comparison.png}
  \caption{QAOA approximation ratio compared to classical baselines for $n = 10$, 3-regular MaxCut. Horizontal lines mark random assignment (AR $= 0.5$), greedy (AR $\approx 0.933$), and Goemans-Williamson SDP (AR $\approx 1.000$, with 0.8786 theoretical floor). QAOA noiseless and noisy curves show mean $\pm$ 95\% CI.}
  \label{fig:baselines}
\end{figure}

\begin{figure}[h]
  \centering
  \includegraphics[width=\linewidth]{figures/fig6_divergence_heatmap.png}
  \caption{Heatmap of mean QAOA approximation ratio across the $(p, n)$ grid. Left: noiseless statevector simulation. Right: FakeSherbrooke noisy simulation. Color scale ranges from AR $= 0.5$ (red) to AR $= 1.0$ (green). Cell values are annotated. The noise-induced erosion of high-$p$ performance is visible for $n \geq 10$.}
  \label{fig:heatmap}
\end{figure}

% ─────────────────────────────────────────────
\section{Discussion}
% ─────────────────────────────────────────────

\subsection{Physical Interpretation of the Divergence Depth}

The divergence depth $p^*$ emerges from a competition between two effects: increasing approximation ratio with $p$ in noiseless simulation, and accumulating gate errors in hardware execution. For MaxCut on a 3-regular graph with $n = 10$, each QAOA layer contains $3n/2 = 15$ logical CNOT gates. After transpilation to IBM's heavy-hex topology, SWAP routing typically multiplies the physical two-qubit gate count by a factor of $\sim 3$, yielding approximately 45 ECR gates per layer. At $p = 3$, this totals $\sim 135$ two-qubit gates. With an ECR error rate of $\varepsilon \approx 0.5\%$, the probability of at least one gate error is approximately $1 - (1 - \varepsilon)^{135} \approx 49\%$. This quantitative argument explains why hardware performance peaks and then declines around $p = 2$--$3$ for $n \geq 10$.

This is consistent with the noise-induced barren plateau theory of Wang et al.\ \cite{wang2021noisebarren}: as $np$ grows, gradient information decays exponentially, making even gradient-free optimization unreliable.

\subsection{Noise Model Fidelity}

Our calibrated depolarizing noise model ($\varepsilon_{2Q} = 0.5\%$, $\varepsilon_{\text{RO}} = 1.0\%$) consistently underestimates the noiseless AR by a gap that grows from $\sim 0.019$ at $p = 1$ to $\sim 0.066$ at $p = 5$ for $n = 10$. Validation against real IBM hardware remains pending; once hardware data is collected, the mean absolute error between model and hardware will quantify whether IBM Eagle R3 calibration data accurately predicts production device noise. Based on comparable studies \cite{pelofske2022quantum}, deviations of $\sim 0.02$--$0.05$ AR units are typical, suggesting our model is in the correct regime. Practitioners can use this noise model to estimate whether a proposed circuit depth is viable before committing IBM Quantum quota.

\subsection{Limitations}

Several limitations constrain the scope of our conclusions. First, we study only $n \leq 12$ qubits. For larger instances, T$_1$/T$_2$ decoherence (rather than gate errors) may become the dominant noise source as circuit execution time grows. Second, we use the fixed-angle protocol: parameters are optimized on the noiseless simulator and transferred to hardware, without re-optimization under noise. Hardware-optimized angles may outperform fixed angles, potentially shifting $p^*$ to higher values. Third, real hardware experiments were not conducted in this study; all noisy results use a calibrated simulation model. Direct hardware validation would strengthen confidence in the noise model fidelity. Fourth, we do not apply error mitigation techniques (zero-noise extrapolation, probabilistic error cancellation), which could extend the effective depth budget.

\subsection{Practitioner Implications}

For practitioners targeting current IBM superconducting hardware (ECR error rate $\sim 0.3$--$0.5\%$), our results suggest the following heuristic: for $n \geq 10$, budget $p \leq 3$ to maintain approximation ratio within 5\% of noiseless simulation. The calibrated depolarizing noise model used here can serve as a planning tool for estimating viable circuit depths before committing hardware quota.

More broadly, our divergence depth framework provides a device-agnostic diagnostic tool. As hardware error rates improve (e.g., IBM Heron R2 ECR error rate $\sim 0.1$--$0.2\%$), $p^*$ should scale inversely with $\varepsilon$, allowing deeper circuits. Measuring $p^*(n, \varepsilon)$ systematically across hardware generations would provide a principled performance roadmap.

% ─────────────────────────────────────────────
\section{Conclusion}
% ─────────────────────────────────────────────

We have introduced the divergence depth $p^*(n)$ as an operationalizable, empirically measurable quantity characterizing the maximum QAOA circuit depth before noise-induced degradation dominates. Benchmarking QAOA for MaxCut across $n \in \{6, 8, 10, 12\}$ qubits and $p \in \{1, 2, 3, 4, 5\}$ layers, with 30 random instances per configuration, we find $p^* = 4$ consistently across all tested system sizes under IBM Eagle R3-representative noise ($\varepsilon_{2Q} = 0.5\%$). Noisy QAOA at its best depth ($p = 3$) achieves AR $= 0.871$ at $n = 10$, remaining $12.9\%$ below the Goemans-Williamson SDP, consistent with no demonstrated practical advantage over polynomial-time classical algorithms at these scales. Validation against real IBM superconducting hardware remains as immediate future work.

Future work should extend this benchmark to larger $n$ using IBM Heron and Flamingo processors, validate the noise model against real hardware, apply error mitigation techniques (zero-noise extrapolation, probabilistic error cancellation) to assess their effect on $p^*$, and explore whether graph families beyond 3-regular MaxCut exhibit different divergence behavior. The complete codebase, raw data, and figure generation scripts are available on GitHub under the MIT license.

% ─────────────────────────────────────────────
\bibliographystyle{unsrt}
\bibliography{bibliography}

\end{document}
